\section{References and Back Matter}

\draftinstr{Back matter is adapted from
\texttt{inputs/phase-1-trial-protocol.zip},
\texttt{sections/sec-12-references-backmatter.tex}, and re-scoped to an advocacy
paper: the abbreviation list keeps only terms this paper actually uses, the
amendment history becomes a version history, and a figure inventory is added
because thirty figures need an index.}

\subsection{Abbreviations}

\vspace{0.30em}

\noindent\begin{tabularx}{\textwidth}{L{2.3cm} Y L{2.3cm} Y}
\toprule
\textbf{Term} & \textbf{Meaning} & \textbf{Term} & \textbf{Meaning} \\
\midrule
AE & Adverse event & IRB & Institutional Review Board \\
CA 19-9 & Carbohydrate antigen 19-9 & ISGPS & International Study Group on Pancreatic Surgery \\
CRO & Contract research organisation & LLM & Large language model \\
DL1 to DL3 & Dose levels 1 to 3 & OS & Overall survival \\
DLT & Dose-limiting toxicity & PDAC & Pancreatic ductal adenocarcinoma \\
DSMB & Data and Safety Monitoring Board & PFS & Progression-free survival \\
ECOG & Eastern Cooperative Oncology Group & R0 & Microscopically clear resection margin \\
e-stop & Emergency stop & RECIST & Response Evaluation Criteria in Solid Tumors \\
FDA & Food and Drug Administration & SAE & Serious adverse event \\
IDE & Investigational Device Exemption & SMV & Superior mesenteric vein \\
IND & Investigational New Drug & USL & Unified Safety Level \\
\bottomrule
\end{tabularx}

\vspace{0.30em}

\draftinstr{Add any abbreviation introduced by \texttt{full-patient} that is not
in this list, and remove any listed term that the finished paper never uses.}

\subsection{Figure inventory}

\draftinstr{Populate all thirty rows from the five stage READMEs. Verify the
numbering is continuous from 1 to 30 with no gaps and no duplicates, and that
the type counts are Mermaid 9, D2 7, PlantUML 5, Graphviz 5, and Diagrams
(Python) 4.}

\vspace{0.30em}

\noindent\begin{tabularx}{\textwidth}{C{1.0cm} L{3.6cm} C{1.2cm} Y}
\toprule
\textbf{Fig} & \textbf{Type} & \textbf{\S} & \textbf{Subject} \\
\midrule
1 & Mermaid-type & 1 & The seven commitments made to the participant \\
2 & Graphviz-type & 1 & Accountability chain rooted at the participant \\
3 & Mermaid-type & 2 & Journey schema with every decision point marked \\
4 & D2-type & 2 & Schedule of Activities from the participant's chair \\
5 & D2-type & 3 & Twenty-one concerns in six disjoint families \\
6 & Graphviz-type & 3 & Every concern wired to its answering clause \\
7 & Mermaid-type & 3 & Concern prevalence against answer completeness \\
8 & PlantUML-type & 3 & Safeguards as capabilities the participant holds \\
9 & Diagrams (Python) & 3 & Where each concern physically lives \\
10 & Mermaid-type & 4 & Objective, endpoint, and the sentence it authorises \\
11 & D2-type & 4 & Endpoint registry with a plain-meaning field \\
12 & PlantUML-type & 5 & Participant status as a guarded state machine \\
13 & Graphviz-type & 5 & The 3+3 rule as an auditable decision tree \\
14 & Mermaid-type & 6 & Eligibility as a two-way gate \\
15 & PlantUML-type & 6 & Direct booking and sponsor response \\
16 & D2-type & 6 & The five choices the participant keeps \\
17 & Diagrams (Python) & 7 & Operating room and on-premises stack \\
18 & PlantUML-type & 7 & Emergency-stop timing budget \\
19 & Mermaid-type & 7 & Advise, approve, execute, and log \\
20 & Graphviz-type & 7 & Hazard fault tree with named barriers \\
21 & D2-type & 7 & Force caps and vascular no-fly envelope \\
22 & PlantUML-type & 8 & Four withdrawal routes and their data outcomes \\
23 & Mermaid-type & 8 & Consent lifecycle with forced re-consent \\
24 & Diagrams (Python) & 9 & Data pipeline with per-store read lists \\
25 & Mermaid-type & 9 & Participant visit calendar with totals \\
26 & D2-type & 9 & Ten robot instruction cards, PDAC-scoped \\
27 & Mermaid-type & 10 & Five-panel quantitative reassurance dashboard \\
28 & Graphviz-type & 10 & Evidence provenance for every quoted number \\
29 & D2-type & 11 & Responsibility matrix, one accountable per row \\
30 & Diagrams (Python) & 12 & Post-trial continuity and the payer per item \\
\bottomrule
\end{tabularx}

\vspace{0.30em}

\subsection{Version history}

\vspace{0.30em}

\noindent\begin{tabularx}{\textwidth}{L{2.0cm} L{2.6cm} Y}
\toprule
\textbf{Version} & \textbf{Date} & \textbf{Change} \\
\midrule
Draft 1.0 & July 31, 2026 & First release. Thirteen sections, thirty figures, twenty-one documented concerns answered, three-stage build from diagram sources through draft and full to final. \\
\bottomrule
\end{tabularx}

\vspace{0.30em}

\bmhead{Data and code availability}

\draftinstr{Roughly 500 characters. Every figure source, every LaTeX source, and
every stage narrative is in the public repository at
\texttt{github.com/kevinkawchak/robotic-surgeries/tree/main/patient-robot-advocacy}.
No participant data exists, because no participant exists: this is an
independent advocacy paper about a protocol, not a report of a conducted study.}

\bmhead{Acknowledgments}

\draftinstr{Roughly 400 characters. Acknowledge the two dated research passes at
\texttt{research/research-a.md} (Gemini 3.1 Pro) and
\texttt{research/research-b.md} (ChatGPT 5.6 Thinking Extended), and state that
the LaTeX build was adapted using Claude Code Opus 5.}

\bmhead{Ethical disclosures}

\draftinstr{Roughly 500 characters. No human participants, no animals, no
identifiable data. This is an independent research paper and practical adoption
guide, not an active IND or IDE, not an approved protocol, and not medical or
regulatory advice. Not endorsed by the FDA, NIH, HHS, an IRB, ICH, or any
sponsor, CRO, site, regulator, or medical society.}

\bmhead{Rights and permissions}

\draftinstr{Roughly 300 characters. Released under CC BY 4.0. Reproduced United
States Government regulatory text is used under 17 U.S.C. \S 105.}

\bmhead{Cite this article}

\draftinstr{Give the citation in the author's usual form, with the DOI as text
and as a clickable link, and with the repository URL. Both the paper (v1.0) and
the repository (v1.0.0) are named.}

\vspace{0.30em}

\noindent Kawchak K. \textit{Patient Robot Advocacy: A Phase 1, First-in-Human,
PDAC Clinical Trial Protocol of a LLM-Directed Robotic Whipple with Daraxonrasib
(RMC-6236)}. Draft 1.0. Zenodo; 2026.
\href{https://doi.org/10.5281/zenodo.xxxxxxxx}{10.5281/zenodo.xxxxxxxx}.
Repository v1.0.0:
\href{https://github.com/kevinkawchak/robotic-surgeries/tree/main/patient-robot-advocacy}{github.com/kevinkawchak/robotic-surgeries/tree/main/patient-robot-advocacy}.

\bmhead{References}

\bibliographystyle{ieeetr}
\bibliography{references}
